\scanpage{74}\mspage{36}
% archive sheet number: 73
\deleted{\uline{Proposition}}\inserted{(9)} Toute arête de $\widetilde P_M$ est contenue dans \deleted{une unique face} \inserted{une unique face}, d'où
\[
 \widetilde A_M\xrightarrow[\text{degré }3]{}\widetilde F_M
\]
et une application composée
\diagram{\begin{tikzcd}[ampersand replacement=\&,column sep=large]
{\text{(10)}} \&\widetilde A_M\arrow[r]\arrow[rr,bend right=22,"\widetilde\tau_A"'] \&\widetilde F_M\arrow[r,"\widetilde\tau_F"] \&\unsure{\pi|\Pi}_M
\end{tikzcd}}
De \unsure{$=$|même}, avec la définition évidente des arêtes orientées, on a
\diagram{\begin{tikzcd}[ampersand replacement=\&,column sep=large]
{\text{(11)}} \&{\underset{\substack{\text{arêtes orientées}\\\text{de }\widetilde A_M}}{\widetilde{\overrightarrow A}_M}}\arrow[r,"\text{degré }3"]\arrow[rr,bend right=31,"\widetilde{\overrightarrow\tau}_A"'] \&\widetilde{\overrightarrow F}_M\arrow[r,"\widetilde{\overrightarrow\tau}_F"] \&\overrightarrow\pi_M
\end{tikzcd}}
\margin{\uline{NB $\widetilde{\overrightarrow\tau}_A(\check{\overrightarrow a})=-\widetilde{\overrightarrow\tau}_A(\overrightarrow a)$}}
% slanted
\margin{(10) \deleted{\unsure{peut}} être vu comme déduit de (11) en passant aux quotients par $\overrightarrow a\longmapsto\check{\overrightarrow a}$}
On définit $\widetilde A_{M,i}$ \deleted{\unsure{en}} (pour $i\in\pi_M$) comme l'image inverse de $\{i\}$ par $\widetilde{\overrightarrow\tau}_A$, d'où
\[
 \widetilde A_M=\coprod_{i\in\pi}\widetilde A_{M,i}\,.\tag*{(12)}
\]
\[
\boxed{\begin{gathered}
 \deleted{\widetilde A_{M,i}\text{ et }\widetilde F_{M,i}\text{ se déterminent mutuellement}}\\
 \left\{\begin{aligned}
 &\deleted{\widetilde A_{M,i}=\{a\in\mathfrak P_2(\widetilde S_M)\mid\exists f\in\widetilde F_{M,i},f\supset a\}}\\
 &\deleted{\widetilde F_{M,i}=\{f\in\mathfrak P_3(\widetilde S_M)\mid\forall a\in\mathfrak P_2(f),a\in\widetilde A_{M,i}\}}
 \end{aligned}\right.
\end{gathered}}\tag*{\deleted{(13)}}
\]
\deleted{\unsure{Dans}} On peut considérer que $(\widetilde F_{M,i})_{i\in\pi_M}$ définit une famille de structures triangulées sur le $=$ ens. de sommets $\widetilde S_M$, paramétré par $\pi_M$.
% slanted
\margin{\deleted{\unsure{ou encore, on a} une structure}}
Il faudra à l'occasion interpréter la structure
% slanted
\margin{\uline{tirer au} \uline{clair}}
\[
 \pi_M\text{ torseur sous }\deleted{\unsure{\Lambda}} Z'_{\Lambda}=\Lambda^*/\{\pm1_\Lambda\},\tag*{(1\deleted{\unsure{5}}3)}
\]
et aussi
\[
 \overrightarrow\pi_M\text{ torseur sous }\Lambda^*\text{ qui ``relève'' }\pi_M\,.\tag*{(1\deleted{\unsure{6}}4)}
\]
